%% =====================================================================
%%  Welcome : la promesse de la méthode, les niveaux, le rythme
%% =====================================================================
\liminaire{bienvenue}{Welcome!}{Everything you need to know before you start}

{\fontsize{11.4}{15.6}\selectfont
You have never written a line of code? Or you tried once, got lost in the jargon and gave up? This
book is for you. It assumes \textbf{no prior knowledge}: no programming, no advanced mathematics,
nothing but the ability to use a computer. Every idea is explained in plain English, shown in
programs that you can run, and practised until it becomes natural.\par\vspace{3mm}

Programming is not reserved for geniuses. It is a skill, like cooking or playing the guitar: you
learn it step by step, by doing. In this first volume, you will write programs that calculate, talk
to their user, make decisions, repeat actions, organise themselves into functions and keep their
data in files. Along the way, you will build five complete projects\,---\,and you will start to
think like a programmer.\par}

\otIntertitre{By the end of this volume, you will be able to\ldots}
\begin{paire}
\begin{tcolorbox}[otBase,colback=otNavyTint,colframe=otNavyTint]
\begin{list}{\otCheck}{\setlength{\leftmargin}{7mm}\setlength{\labelwidth}{4mm}\setlength{\labelsep}{3mm}
  \setlength{\itemsep}{1.6mm}\setlength{\topsep}{0pt}\setlength{\parsep}{0pt}}\fontsize{10.4}{13}\selectfont
\item install Python and work with IDLE, its editor;
\item calculate, and keep values in variables;
\item work with text: build it, cut it, transform it;
\item write programs that ask questions and react to the answers;
\end{list}
\end{tcolorbox}
\begin{tcolorbox}[otBase,colback=otNavyTint,colframe=otNavyTint]
\begin{list}{\otCheck}{\setlength{\leftmargin}{7mm}\setlength{\labelwidth}{4mm}\setlength{\labelsep}{3mm}
  \setlength{\itemsep}{1.6mm}\setlength{\topsep}{0pt}\setlength{\parsep}{0pt}}\fontsize{10.4}{13}\selectfont
\item repeat actions with loops and create simple games;
\item organise your code into functions that you can reuse;
\item handle errors so that your programs do not crash;
\item store data in lists, dictionaries and files.
\end{list}
\end{tcolorbox}
\end{paire}

\otIntertitre{Level 1, the first step of a five-volume journey}
The collection takes you from your very first instruction to the skills of a professional
developer. This volume is the foundation: everything in the following volumes builds on it.\par\vspace{4mm}

\noindent\begin{tikzpicture}
  \foreach \n/\nom/\t/\c in {1/Beginner/Volume 1/otGreen,2/Elementary/Volume 2/otGreenDark,
      3/Intermediate/Volume 3/otNavy!70,4/Advanced/Volume 4/otNavy,5/Expert/Volume 5/otNavyDark}{%
    \pgfmathsetmacro{\x}{(\n-1)*34.2}
    \pgfmathsetmacro{\h}{10+(\n-1)*7.5}
    \fill[\c,rounded corners=1.2mm] (\x mm,0) rectangle ++(32mm,\h mm);
    \node[text=white,font=\otBold\fontsize{15}{16}\selectfont] at (\x mm+16mm,\h mm-5.6mm) {\n};
    \node[anchor=north,font=\otSemi\fontsize{8.4}{10}\selectfont,text=otInk] at (\x mm+16mm,-1.6mm) {\nom};
    \node[anchor=north,font=\otRegular\fontsize{7.8}{9.4}\selectfont,text=otMuted] at (\x mm+16mm,-5.6mm) {\t};}
  \node[anchor=south,fill=otOrange,text=white,rounded corners=1.4mm,inner xsep=2.4mm,inner ysep=1.2mm,
    font=\otSemi\fontsize{8}{9}\selectfont] (ici) at (16mm,17mm) {You are here};
  \draw[otOrange,line width=1pt,-{Stealth[length=2mm]}] (ici.south) -- (16mm,11mm);
\end{tikzpicture}

\clearpage
\otIntertitre{Your rhythm: one chapter a week}
A chapter takes about \textbf{three hours} of work. The most effective way is to spread them over
\textbf{four sessions of about 45 minutes}, listed on the first page of each chapter, and to keep a
little time to replay the code. At this pace, you will finish the volume in about six months,
projects and checkpoints included.\par\vspace{4mm}

\noindent\begin{tikzpicture}
  \foreach \j/\d/\s/\c in {1/Mon/Session 1/otNavy,2/Tue/Session 2/otViolet,3/Wed/Replay the code/otMuted,
      4/Thu/Session 3/otTeal,5/Fri/Session 4/otGreenDark,6/Sat/Quick quiz/otOrangeDark,7/Sun/Rest or free coding/otMuted}{%
    \pgfmathsetmacro{\x}{(\j-1)*24.4}
    \fill[\c!12,rounded corners=1.4mm] (\x mm,0) rectangle ++(22.6mm,20mm);
    \fill[\c,rounded corners=1.4mm] (\x mm,13mm) rectangle ++(22.6mm,7mm);
    \node[text=white,font=\otSemi\fontsize{8.4}{10}\selectfont] at (\x mm+11.3mm,16.5mm) {\d};
    \node[text width=20mm,align=center,font=\otMedium\fontsize{7.8}{9.4}\selectfont,text=otInk] at (\x mm+11.3mm,6.5mm) {\s};}
\end{tikzpicture}

\otIntertitre{Six keys to success}
\begin{tcbraster}[raster columns=3,raster equal height=rows,raster column skip=3.5mm,raster row skip=3.5mm,
  raster before skip=1mm,raster after skip=2mm]
\otCartePrg{clavier}{otNavy}{Type the code yourself}{Copying and pasting teaches you nothing. Typing every line, even the short ones, trains your eyes and your fingers.}
\otCartePrg{lecture}{otGreenDark}{Run it often}{Press F5 after every small change. The sooner you see the result, the sooner you understand.}
\otCartePrg{bug}{otOrangeDark}{Break things on purpose}{Change a value, delete a parenthesis, and watch what happens. Mistakes are excellent teachers.}
\otCartePrg{oeil}{otViolet}{Read the messages}{Python's error messages tell you what went wrong and where. Read them from the bottom up.}
\otCartePrg{bulle}{otTeal}{Explain it out loud}{If you can explain a program line by line to someone else, you have really understood it.}
\otCartePrg{horloge}{otRedDark}{A little every day}{Forty-five minutes a day beat five hours on Sunday: regularity is everything.}
\end{tcbraster}

\begin{astuce}[Before you begin]
Read the next pages: \textbf{How to use this book} (p.~\pageref{lim-mode}) and \textbf{Getting the
code} (p.~\pageref{lim-code}). Ten minutes well spent: you will know exactly how to work with every
page.
\end{astuce}
